\chapter{The Rosetta Stone: Types and Expressions}
\label{ch:rosetta}

\chapabstract{The first half of the reference core. For every \VHDL{} type,
literal, operator and expression rule you have used during your degree, this
chapter gives the \SV{} form beside it, says what is identical, what is only
superficially similar, and what has no counterpart at all. The centre of it is
the width and signedness rules, which are the ones that silently change the
arithmetic your design performs. The behavioural half, processes and
statements and everything else that happens over time, is \cref{ch:rosetta2}.}

A translation dictionary is useful only if it also tells you where the two
languages disagree. Most of \VHDL{} has a direct \SV{} equivalent, and for
those constructs a side-by-side pane is enough. A small but dangerous
minority looks equivalent and is not: the width rules of arithmetic, the
meaning of an unsized literal, the resolution of a multiply driven signal.
Those get a trap box, and you should read every one of them.

The dictionary is split across two chapters because it divides cleanly in two,
and because reading four hundred mappings in one sitting is not how anybody
uses a reference. This chapter is the \emph{static} half: how a design is
broken into files, what a port and a parameter are, what the types are, what a
literal means, and what happens to widths and signedness when you write an
expression. Nothing here depends on time passing. \Cref{ch:rosetta2} is the
half that does: processes, sequential statements, \code{generate},
instantiation, subprograms, packages, attributes, reporting and delays. It
also carries the complete mapping table for both halves, so if you are looking
for the page to photocopy, it is at the end of the next chapter and not this
one.

The layout is constant. Each section opens with one or more \emph{rosetta}
blocks, \VHDL{} on the left, \SV{} on the right. The prose after the block
answers three questions in order: what is the same, what is different, what is
missing. Where a construct exists in only one language the missing pane is
filled with the closest idiom and the text says plainly that it is not a
translation.

Throughout, \VHDL{} means IEEE 1076-2008 and \SV{} means IEEE 1800-2017. Where
a \VHDL{}-2019 feature would change the answer, the text says so and warns you
that tool support is thin.

% ===========================================================================
\section{Design units, files and the compilation model}
\label{sec:rosetta:units}
% ===========================================================================

\VHDL{} has a strict notion of a design unit and an equally strict analysis
order. \SV{} inherited Verilog's file-oriented, preprocessor-driven model,
which is looser in every respect. This is the first thing to internalise,
because it explains a great deal of what follows.

\index{design unit}\index{compilation unit}

\subsection{Entity and architecture versus module}

\begin{rosetta}
\begin{rleft}
library ieee;
use ieee.std_logic_1164.all;
use ieee.numeric_std.all;

entity counter is
  generic (W : positive := 8);
  port (
    clk : in  std_logic;
    rst : in  std_logic;
    en  : in  std_logic;
    q   : out std_logic_vector(W-1 downto 0)
  );
end entity counter;

architecture rtl of counter is
  signal cnt : unsigned(W-1 downto 0);
begin
  process (clk) is
  begin
    if rising_edge(clk) then
      if rst = '1' then
        cnt <= (others => '0');
      elsif en = '1' then
        cnt <= cnt + 1;
      end if;
    end if;
  end process;
  q <= std_logic_vector(cnt);
end architecture rtl;
\end{rleft}
\begin{rright}
`default_nettype none

module counter #(
  parameter int unsigned W = 8
) (
  input  wire              clk,
  input  wire              rst,
  input  wire              en,
  output logic [W-1:0]     q
);

  logic [W-1:0] cnt;

  always_ff @(posedge clk) begin
    if (rst)      cnt <= '0;
    else if (en)  cnt <= cnt + 1;
  end

  assign q = cnt;

endmodule

`default_nettype wire
\end{rright}
\end{rosetta}

The two files describe the same hardware, and almost every line has a
counterpart. The differences that matter are structural.

\VHDL{} splits the black box in two. The \kw{entity} is the interface, the
\kw{architecture} is one possible implementation of it, and the pair is bound
either implicitly (last analysed architecture wins) or explicitly by a
\kw{configuration}. \SV{} has a single \kw{module}\index{module} that carries
both. There is no separate interface unit to compile, no analysis order
between them, and no default binding rule to remember.

\SV{} does have a \kw{config} construct, kept from Verilog-2001, that plays
roughly the role of a \VHDL{} configuration by selecting which library cell
satisfies which instance. In practice nobody writes them for \RTL{}, tool
support varies, and the same job is done with library search order on the
command line. Treat configurations as a \VHDL{} idiom without a practical
\SV{} equivalent.

\subsection{Several architectures for one entity}

The \VHDL{} habit of writing \code{rtl}, \code{behavioural} and
\code{structural} architectures for one entity has no direct translation. \SV{}
gives you three replacements, in decreasing order of preference.

\begin{rosetta}
\begin{rleft}
architecture rtl of alu is
begin
  -- synthesisable description
end architecture rtl;

architecture behav of alu is
begin
  -- untimed reference model
end architecture behav;

-- selected by a configuration,
-- or by analysis order
\end{rleft}
\begin{rright}
module alu #(
  parameter bit BEHAVIOURAL = 0
) ( /* ports */ );
  if (BEHAVIOURAL) begin : g_ref
    // untimed reference model
  end else begin : g_rtl
    // synthesisable description
  end
endmodule

// or: two modules, alu_rtl and
// alu_behav, chosen by the
// instantiating code
\end{rright}
\end{rosetta}

The first replacement, shown above, is a parameter tested inside a
\code{generate} region. It keeps one module name and one port list, which is
exactly what a \VHDL{} entity gave you. The second is a preprocessor
\chTwoTick{ifdef}, which works but hides the choice from the elaborator. The
third, and the most common in industry, is simply two modules with different
names and a wrapper that instantiates one of them. \SV{} has no notion of
``the same interface, different body'', so the interface is repeated.

\begin{notebox}{Why the split existed}
The entity/architecture split is a 1980s design decision aimed at supporting
separate compilation and technology retargeting: analyse the entity once,
supply a gate-level architecture later. Verilog solved the same problem with
its library and \code{config} mechanism, which nobody used, and in practice with
file lists. The industry converged on file lists. This is a case where the
weaker language mechanism won because the tooling around it was simpler.
\end{notebox}

\subsection{Libraries, use clauses and imports}

\begin{rosetta}
\begin{rleft}
library ieee;
use ieee.std_logic_1164.all;
use ieee.numeric_std.all;

library work;
use work.cordic_pkg.all;
use work.cordic_pkg.data_t;
\end{rleft}
\begin{rright}
// no library clause needed:
// std types are built in

import cordic_pkg::*;
import cordic_pkg::data_t;

`include "cordic_defs.svh"
\end{rright}
\end{rosetta}

\index{library!VHDL}\index{import}

There is no \kw{library} clause in \SV{}. A \VHDL{} library is a named
container of analysed design units, created by the tool
(\code{vcom -work mylib}), and referenced by name from the source. \SV{}
libraries exist in the \LRM{} but the everyday model is: everything you compile
lands in one namespace, and modules are found by name.

\code{use work.pkg.all} maps to \code{import pkg::*}. The two are close enough
that the translation is mechanical, but the resolution rules differ, and
\cref{sec:rosetta:packages} treats them properly. The essential difference:
\code{use ... all} makes the names visible immediately, while
\code{import pkg::*} makes them visible \emph{only on reference and only if
nothing else already declares that name}. A wildcard import that would collide
is not an error until you actually use the ambiguous name.

The two lines that have no \VHDL{} counterpart at all are the absence of
\kw{library} and the presence of \chTwoTick{include}. The second is a textual
file inclusion performed by a preprocessor before the language proper sees the
text. \VHDL{} has nothing like it, and never needed it, because a package is a
first-class compiled unit.

\subsection{The compilation unit and \code{\$unit}}

In \VHDL{} every design unit belongs to a library and the analysis order is
enforced: a package body cannot be analysed before its declaration, an
architecture cannot be analysed before its entity. Change a package and every
unit that used it is marked out of date. This is a genuine strength and \SV{}
has nothing equivalent for the general case.

\SV{} defines a \emph{compilation unit}\index{compilation unit}: the set of
source files compiled together in one invocation of the tool. Declarations
made outside any module, package or interface live in the compilation-unit
scope, addressed as \code{\$unit}.

\begin{svcode}[caption={A compilation-unit-scope declaration. Legal, and a bad idea.},label={lst:rosetta:dollarunit}]
// file: types.sv -- outside any module or package
typedef logic [15:0] word_t;      // lands in $unit
localparam int unsigned NREG = 32;

module regfile (input word_t d, output word_t q);
  word_t mem [NREG];              // visible here only if compiled together
endmodule
\end{svcode}

The trap is that ``compiled together'' is defined by the tool invocation, not by
the language. \tool{QuestaSim} with two \code{vlog} commands gives you two
compilation units and \code{word\_t} is invisible in the second.
\tool{Verilator}, \tool{VCS} and \tool{Xcelium} all have their own defaults and
their own switches. A design that relies on \code{\$unit} compiles for one
person and not for the next.

\begin{ruleofthumb}
Put nothing in \code{\$unit}. Every type, constant and function belongs in a
package; every macro belongs in an included header.
\end{ruleofthumb}

\subsection{The preprocessor}

Verilog's preprocessor is a separate, weaker language that runs over the text
before compilation. It is the only place in \SV{} where the backtick appears.

\begin{rosetta}
\begin{rleft}
-- VHDL has no preprocessor.
-- The nearest equivalents:

constant DEPTH : positive := 16;

-- conditional analysis, VHDL-2019
-- (thin tool support):
--   if TOOL_TYPE = "SIMULATION" generate
\end{rleft}
\begin{rright}
`timescale 1ns/1ps
`define DEPTH 16
`define MAX(a,b) (((a)>(b))?(a):(b))

`ifdef SYNTHESIS
  // synthesis-only text
`elsif FORMAL
  // formal-only text
`else
  // simulation-only text
`endif

`undef DEPTH
`include "defs.svh"
\end{rright}
\end{rosetta}

\index{preprocessor}\index{macro!Verilog}

Macros are textual, untyped and unscoped. \chTwoTick{DEPTH} is not a constant:
it is the characters \code{16} pasted in. A macro defined in one file is
visible in every file compiled after it in the same compilation unit, which
means macro definitions leak across your whole design in file order. Prefer a
package \code{localparam} for every value, and reserve macros for the two jobs
they alone can do: conditional inclusion of text, and generating repetitive
declarations that the language cannot parameterise.

\chTwoTick{timescale} deserves its own trap and gets one in
\cref{sec:rosetta:time}.

\begin{tipbox}{Always start an RTL file with \chTwoTick{default\_nettype none}}
Verilog creates a one-bit wire for any undeclared identifier used in a port
connection or on the right of a continuous assignment. A misspelt signal name
therefore compiles cleanly and produces a silently disconnected net. Putting
\chTwoTick{default\_nettype none} at the top of every file turns that into an
error. Restore it with \chTwoTick{default\_nettype wire} at the bottom of the
file so that you do not break code compiled after yours.
\end{tipbox}

\begin{gotcha}{Implicit nets: the typo that compiles}
\VHDL{} rejects an undeclared signal. Verilog invents one.
\begin{svcode}
module top (input wire a, b, output logic y);
  wire and_out;
  assign and_out = a & b;
  assign y       = and_ou;    // typo. Compiles. y is z forever.
endmodule
\end{svcode}
The name \code{and\_ou} is undeclared, so Verilog creates a new one-bit net that
nothing drives. Simulation shows \code{z}, synthesis shows a constant or an
optimised-away output, and there is no message. \chTwoTick{default\_nettype none}
is the only reliable defence, and it is not the default because thirty years of
existing code depends on the old behaviour.
\end{gotcha}

% ===========================================================================
\section{Ports}
\label{sec:rosetta:ports}
% ===========================================================================

\begin{rosetta}
\begin{rleft}
entity fifo is
  port (
    clk   : in  std_logic;
    rst_n : in  std_logic;
    wr_en : in  std_logic;
    din   : in  std_logic_vector(7 downto 0);
    rd_en : in  std_logic;
    dout  : out std_logic_vector(7 downto 0);
    full  : out std_logic;
    empty : buffer std_logic;
    sda   : inout  std_logic
  );
end entity fifo;
\end{rleft}
\begin{rright}
module fifo (
  input  wire        clk,
  input  wire        rst_n,
  input  wire        wr_en,
  input  wire  [7:0] din,
  input  wire        rd_en,
  output logic [7:0] dout,
  output logic       full,
  output logic       empty,
  inout  wire        sda
);
endmodule
\end{rright}
\end{rosetta}

\index{port!direction}

\subsection{Modes and directions}

\cref{tab:rosetta:portmodes} gives the mapping. Of the four modes you will
actually meet, three translate directly; the fifth entry in the table,
\code{linkage}, is obsolete in \VHDL{} as well.

\begin{table}[htbp]
  \centering
  \caption{Port modes.}
  \label{tab:rosetta:portmodes}
  \small
  \begin{tabular}{@{}lll@{}}
    \toprule
    \VHDL{} mode & \SV{} direction & Notes \\
    \midrule
    \code{in}     & \code{input}  & Read only in both. \\
    \code{out}    & \code{output} & Readable inside the module in \SV{}. \\
    \code{inout}  & \code{inout}  & Must be a net type in \SV{}. \\
    \code{buffer} & \code{output} & No equivalent needed; \code{output} already reads back. \\
    \code{linkage}& ---           & Obsolete in \VHDL{} too. Ignore. \\
    \bottomrule
  \end{tabular}
\end{table}

The interesting entry is \code{buffer}. Before \VHDL{}-2008 an \code{out} port
could not be read inside the architecture, so you either declared an internal
signal and assigned it to the port, or used mode \code{buffer} and accepted its
restrictions on how the port could be connected. \VHDL{}-2008 removed the
restriction: an \code{out} port is now readable. \SV{} never had the
restriction. An \code{output logic} is an ordinary variable that happens to be
visible outside, and you read it as freely as any internal signal:

\begin{svcode}[caption={Reading an output port is normal in \SV{}.},label={lst:rosetta:readout}]
module gray_counter (input wire clk, output logic [3:0] cnt);
  always_ff @(posedge clk)
    cnt <= cnt + 1;            // cnt is read and written; no helper signal
endmodule
\end{svcode}

The habit of declaring \code{signal cnt\_int} and writing
\code{cnt <= cnt\_int} at the end of the architecture is a pre-2008 \VHDL{}
reflex. Do not carry it into \SV{}.

\subsection{Port types: \kw{wire} or \kw{logic}}

A \VHDL{} port has a type and nothing else. A \SV{} port has a direction, a
\emph{data type}, and separately a \emph{net or variable kind}. The default,
if you write nothing, is \kw{wire}, which is a four-state net type resolved
across all its drivers.

The convention used throughout this book, and the one you should adopt, is:

\begin{itemize}
  \item \code{input wire} for inputs. An input is always driven from outside,
        so a net is correct, and it also documents that the module does not
        drive it.
  \item \code{output logic} for outputs. A variable can be assigned from an
        \code{always} block or from a continuous assignment, which covers every
        \RTL{} case.
  \item \code{inout wire} for true bidirectional pads. A variable cannot be
        used, because more than one driver is the whole point.
\end{itemize}

Writing \code{output wire} is legal and forces you to drive the port with a
continuous \code{assign}; writing \code{input logic} is legal and mostly
harmless. Consistency matters more than the choice, but \code{output logic} is
the one that prevents the most mistakes.

\subsection{Default values on ports}

\begin{rosetta}
\begin{rleft}
entity dsp is
  port (
    clk  : in std_logic;
    -- default used if the port is
    -- left unconnected at the
    -- instantiation:
    ce   : in std_logic := '1';
    mode : in std_logic_vector(1 downto 0)
            := "00"
  );
end entity dsp;
\end{rleft}
\begin{rright}
module dsp (
  input wire       clk,
  input wire       ce,
  input wire [1:0] mode
);
  // SystemVerilog has no default
  // port values. An unconnected
  // input is high impedance, and
  // reads as x through any logic.
endmodule
\end{rright}
\end{rosetta}

This is a real loss. \VHDL{} lets you add an optional port to an existing
entity and leave every existing instantiation untouched. \SV{} does not: an
unconnected input port is undriven, and \code{z} propagating into combinational
logic gives \code{x}. The two available workarounds are a parameter instead of
a port, when the value is elaboration-time constant, or an explicit tie-off at
every instantiation site.

\begin{gotcha}{An unconnected input is not zero}
Leaving an input off a \SV{} instantiation is legal and silent. The port is
\code{z}. It does not default to \code{0}, it does not default to the last
value, and most simulators will not warn you. \tool{Verilator} does warn
(\code{PINMISSING}) and this is one of the reasons to run it over your code even
if you simulate elsewhere. Connect every port, every time.
\end{gotcha}

\subsection{ANSI and non-ANSI port lists}

You will meet legacy code that declares ports twice. Recognise it, and never
write it.

\begin{rosetta}
\begin{rleft}
-- VHDL has only one form.
entity adder is
  port (
    a, b : in  std_logic_vector(7 downto 0);
    s    : out std_logic_vector(8 downto 0)
  );
end entity adder;
\end{rleft}
\begin{rright}
// non-ANSI (Verilog-1995). Legacy.
module adder (a, b, s);
  input  [7:0] a;
  input  [7:0] b;
  output [8:0] s;
  reg    [8:0] s;
  // ...
endmodule
\end{rright}
\end{rosetta}

The non-ANSI form names the ports in the header and gives their direction, size
and kind in separate declarations in the body. It is verbose, it lets the two
halves disagree, and it makes parameterised widths awkward. The ANSI form,
introduced in Verilog-2001 and used everywhere in this book, puts everything in
the header. Use it exclusively.

% ===========================================================================
\section{Generics and parameters}
\label{sec:rosetta:params}
% ===========================================================================

\begin{rosetta}
\begin{rleft}
entity fifo is
  generic (
    WIDTH : positive := 8;
    DEPTH : positive := 16;
    NAME  : string   := "fifo"
  );
  port (...);
end entity fifo;

architecture rtl of fifo is
  constant ADDR_W : natural :=
    integer(ceil(log2(real(DEPTH))));
  type mem_t is array (0 to DEPTH-1) of
    std_logic_vector(WIDTH-1 downto 0);
  signal mem : mem_t;
begin
end architecture rtl;
\end{rleft}
\begin{rright}
module fifo #(
  parameter int unsigned WIDTH = 8,
  parameter int unsigned DEPTH = 16,
  parameter string       NAME  = "fifo"
) (...);

  localparam int unsigned ADDR_W =
    $clog2(DEPTH);

  logic [WIDTH-1:0] mem [DEPTH];

endmodule
\end{rright}
\end{rosetta}

\index{parameter}\index{localparam}\index{generic}

\subsection{The mapping}

A \VHDL{} \kw{generic} is an elaboration-time constant supplied by the
instantiating code. So is a \SV{} \kw{parameter}. A \VHDL{} \kw{constant}
declared in the architecture is computed from the generics and cannot be
overridden. So is a \SV{} \kw{localparam}. The correspondence is exact, and it
is worth stating the rule as a discipline:

\begin{ruleofthumb}
\kw{parameter} for anything the instantiator may set, \kw{localparam} for
anything derived from it. Never leave a derived value as a \kw{parameter}: it
can then be overridden inconsistently with the value it was derived from.
\end{ruleofthumb}

\subsection{Types on parameters}

A \VHDL{} generic is always typed. A Verilog parameter historically was not:
its type came from the type of its default value, which was usually an
untyped integer. \SV{} allows and encourages an explicit type, and you should
always write one.

\begin{svcode}[caption={Typed and untyped parameters.},label={lst:rosetta:paramtypes}]
parameter                WIDTH_A = 8;      // type taken from 8: 32-bit signed int
parameter int unsigned   WIDTH_B = 8;      // explicit, checked at override
parameter logic [7:0]    INIT    = 8'hA5;  // explicit vector
parameter real           FREQ_HZ = 100.0e6;
parameter string         NAME    = "u_fifo";
parameter type           payload_t = logic [31:0];   // a type parameter
\end{svcode}

The last line has no \VHDL{}-2008 equivalent in common use. \SV{} lets a
parameter be a \emph{type}, so a module can be generic over the type of its
payload:

\begin{rosetta}
\begin{rleft}
-- VHDL-2008 generic types.
-- Legal, but tool support is thin
-- and few designers use them.
entity reg_stage is
  generic (type payload_t);
  port (d : in  payload_t;
        q : out payload_t);
end entity reg_stage;
\end{rleft}
\begin{rright}
module reg_stage #(
  parameter type payload_t = logic
) (
  input  wire       clk,
  input  payload_t  d,
  output payload_t  q
);
  always_ff @(posedge clk) q <= d;
endmodule
\end{rright}
\end{rosetta}

\VHDL{}-2008 does have generic types, and generic packages and generic
subprograms with them. Support in the mainstream simulators is partial and in
the synthesis tools it is worse, so most \VHDL{} designers have never used
them. In \SV{} a type parameter is ordinary, well supported, and the standard
way to write a reusable pipeline register, FIFO or arbiter.

\subsection{Overriding at instantiation}

\begin{rosetta}
\begin{rleft}
u_fifo : entity work.fifo
  generic map (
    WIDTH => 16,
    DEPTH => 32
  )
  port map (
    clk => clk,
    ...
  );

-- positional generic map, legal
-- but discouraged:
u_bad : entity work.fifo
  generic map (16, 32)
  port map (...);
\end{rleft}
\begin{rright}
fifo #(
  .WIDTH (16),
  .DEPTH (32)
) u_fifo (
  .clk (clk),
  ...
);

// positional, legal but
// discouraged:
fifo #(16, 32) u_bad (...);

// defparam: obsolete, never use
// defparam u_fifo.WIDTH = 16;
\end{rright}
\end{rosetta}

Named association works in both languages and is the only form you should
write. Note the position of the instance name: \VHDL{} puts the label first,
\SV{} puts it after the parameter list and before the port list. This trips up
every \VHDL{} designer for about a week.

\code{defparam}\index{defparam} is a Verilog-1995 mechanism for reaching into a
hierarchy and changing a parameter of an already-instantiated module from
outside. It is deprecated in IEEE 1800, it defeats separate elaboration, and
several tools reject it. There is no \VHDL{} equivalent because \VHDL{} never
made the mistake. Never write it.

\begin{notebox}{Parameters that must not be overridden}
A \VHDL{} generic can be constrained by its subtype: \code{generic (N : positive)}
rejects zero and negative values at elaboration. \SV{} has no subtype
constraints, so \code{parameter int unsigned N = 8} will happily accept
\code{-1} reinterpreted as a very large unsigned number. Guard it yourself with
an elaboration-time assertion:
\begin{svcode}
if (N < 1 || N > 64) $error("N out of range: %0d", N);
\end{svcode}
placed in a \code{generate} region, or with the \code{\$fatal} form inside an
\code{initial} block. The first is checked at elaboration and is the better
choice.
\end{notebox}

% ===========================================================================
\section{Data types}
\label{sec:rosetta:types}
% ===========================================================================

This is the section where the two languages diverge most, and where a careless
translation costs you a week of debugging. \VHDL{} is strongly typed: the type
system is the main tool the language gives you for catching mistakes. \SV{} is
weakly typed by inheritance from Verilog, which had exactly one type, the bit
vector, and grew everything else on top of it. Almost every expression that
\VHDL{} would reject, \SV{} accepts by converting something silently.

\subsection{The four-state and nine-state logic types}

\begin{rosetta}
\begin{rleft}
-- from ieee.std_logic_1164
signal a : std_ulogic;   -- unresolved
signal b : std_logic;    -- resolved

-- the nine values:
--  'U' uninitialised
--  'X' forcing unknown
--  '0' forcing zero
--  '1' forcing one
--  'Z' high impedance
--  'W' weak unknown
--  'L' weak zero
--  'H' weak one
--  '-' don't care
\end{rleft}
\begin{rright}
logic a;   // 4-state variable
wire  b;   // 4-state net, resolved
bit   c;   // 2-state, 0 or 1 only
reg   d;   // same type as logic

// the four values:
//  0
//  1
//  x  unknown
//  z  high impedance

// no 'U', no 'W'/'L'/'H',
// no '-'
\end{rright}
\end{rosetta}

\index{logic type}\index{std\_logic}\index{four-state}

Take the value systems first. \code{std\_logic} has nine values because it was
designed to model both drive strength and the difference between ``never
assigned'' and ``assigned a conflicting value''. \SV{} has four, because Verilog
modelled strength separately, through net strengths, which nobody uses in
\RTL{}. The mapping is:

\begin{table}[htbp]
  \centering
  \caption{Value systems.}
  \label{tab:rosetta:values}
  \small
  \begin{tabular}{@{}ll>{\raggedright\arraybackslash}p{0.60\linewidth}@{}}
    \toprule
    \code{std\_logic} & \SV{} & Comment \\
    \midrule
    \code{'U'} & --- & Reset value of a \SV{} four-state variable is \code{x}, not a separate value. \\
    \code{'X'} & \code{x} & Direct equivalent. \\
    \code{'0'}, \code{'1'} & \code{0}, \code{1} & Direct equivalent. \\
    \code{'Z'} & \code{z} & Direct equivalent. \\
    \code{'W'}, \code{'L'}, \code{'H'} & --- & Drive strengths; \SV{} keeps these in the (unused) strength system. \\
    \code{'-'} & --- & No value. Don't-care is expressed by \code{casez}/\code{case inside}, not by a value. \\
    \bottomrule
  \end{tabular}
\end{table}

The loss of \code{'U'} matters more than it looks. In \VHDL{} an unassigned
signal starts at \code{'U'} and any logic that touches it produces \code{'U'},
so a missing reset shows up immediately as a distinct colour in the waveform.
In \SV{} an unassigned four-state variable starts at \code{x}, which is the
same value a genuine conflict produces. You lose the ability to distinguish
``never driven'' from ``driven to an unknown''. Two-state types (\kw{bit},
\kw{int}) are worse still: they start at \code{0}, so a missing reset simply
looks like a working reset.

\begin{gotcha}{Two-state types hide missing resets}
\begin{svcode}
bit   [7:0] count_a;    // starts at 0
logic [7:0] count_b;    // starts at x
always_ff @(posedge clk) begin
  count_a <= count_a + 1;   // no reset. Simulates perfectly from 0.
  count_b <= count_b + 1;   // no reset. x forever. Bug visible.
end
\end{svcode}
Silicon does not start at zero. Use four-state types (\kw{logic}) for
everything in the \RTL{}, and reserve \kw{bit}, \kw{int} and friends for
testbench bookkeeping where a two-state value is what you actually mean. The
speed advantage of two-state simulation is real but it is not yours to take at
the \RTL{} level.
\end{gotcha}

\subsection{Resolution and multiple drivers}

\index{resolution function}\index{multiple drivers}

\code{std\_logic} is a \emph{resolved subtype} of \code{std\_ulogic}. The
resolution is done by an ordinary \VHDL{} function, \code{resolved}, written in
\code{std\_logic\_1164} and visible in the source. You can write your own
resolution function for your own type. \code{std\_ulogic} is unresolved, and
two drivers on a \code{std\_ulogic} signal are an elaboration error.

\SV{} moves the choice into the type kind rather than the type.

\begin{rosetta}
\begin{rleft}
-- unresolved: two drivers = error
signal s : std_ulogic;

-- resolved: two drivers merge
-- through function `resolved'
signal r : std_logic;

-- your own resolution:
function my_res (v : my_vec)
  return my_type;
subtype my_res_t is
  my_res my_type;
\end{rleft}
\begin{rright}
// variable: one driver only
logic s;

// net: drivers merge through the
// built-in wire resolution table
wire r;

// user-defined resolution:
//   nettype with a resolution
//   function -- exists, but is
//   not supported for synthesis
\end{rright}
\end{rosetta}

A \kw{logic} is a variable. It may be written by exactly one continuous
assignment, or by procedural assignments from one \code{always} block, and any
second source is an elaboration error in a conforming tool. A \kw{wire} is a
net. Two drivers on a \kw{wire} are legal and are resolved by the built-in wire
resolution: \code{0} against \code{1} gives \code{x}, anything against
\code{z} gives the anything. That table is fixed by the \LRM{} and you cannot
change it.

\SV{}-2012 added user-defined nettypes (\code{nettype real\_net\_t with sum;}),
which is the closest thing to a \VHDL{} resolution function. It is intended for
analogue and mixed-signal modelling, it has no synthesis meaning, and you will
not need it.

\begin{ruleofthumb}
\kw{logic} unless you need multiple drivers. \kw{wire} only for a real
tri-state bus or for module inputs.
\end{ruleofthumb}

\subsection{Vectors, ranges and direction}

\begin{rosetta}
\begin{rleft}
signal a : std_logic_vector(7 downto 0);
signal b : std_logic_vector(0 to 7);
signal u : unsigned(15 downto 0);
signal s : signed(15 downto 0);

-- attributes
a'range          -- 7 downto 0
a'length         -- 8
a'high           -- 7
a'low            -- 0
a'left           -- 7
a'right          -- 0
a'ascending      -- false
a'reverse_range  -- 0 to 7
\end{rleft}
\begin{rright}
logic        [7:0] a;
logic        [0:7] b;
logic       [15:0] u;   // unsigned
logic signed[15:0] s;

// system functions
$bits(a)     // 8
$size(a)     // 8
$high(a)     // 7
$low(a)      // 0
$left(a)     // 7
$right(a)    // 0
$increment(a)// 1 (i.e. [7:0])
// no $range, no reverse_range
\end{rright}
\end{rosetta}

\index{downto}\index{vector range}

Both languages let you write the range in either direction. Both make the
descending form the sane choice for a number, because bit $n$ then has weight
$2^n$. The difference is one of convention strength: in \VHDL{} you meet
\code{0 to N-1} vectors regularly, often in bus standards translated from a
document; in \SV{} an ascending packed range is so rare that many engineers
have never seen one, and some tools handle it poorly.

\begin{ruleofthumb}
Write \code{[N-1:0]} for every packed vector. If you inherit \code{[0:N-1]}
code, convert it before you extend it.
\end{ruleofthumb}

The attribute-to-system-function mapping is nearly one to one, with two
exceptions. \code{'range} has no \SV{} equivalent, because \SV{} has no range
objects; a \VHDL{} \code{for i in a'range loop} becomes
\code{for (int i = 0; i < \$size(a); i++)} or, for unpacked arrays,
\code{foreach (a[i])}. \code{'reverse\_range} has no equivalent either. In the
other direction \code{\$bits} is more general than \code{'length}: it returns
the total number of bits of \emph{any} type, including a packed struct or an
enumeration, which \VHDL{} can only do by hand.

\subsection{The scalar types}

\cref{tab:rosetta:scalars} is the table to memorise. Note in particular that
\SV{} has both \kw{int} and \kw{integer}, that they differ, and that the one
whose name matches \VHDL{}'s is the one you will almost never use.

\begin{table}[htbp]
  \centering
  \caption{Scalar types. ``4-state'' means the type can hold \code{x} and
           \code{z}.}
  \label{tab:rosetta:scalars}
  \small
  \begin{tabular}{@{}llll@{}}
    \toprule
    \VHDL{} & \SV{} & Width and sign & States \\
    \midrule
    \code{bit}            & \code{bit}       & 1, unsigned  & 2 \\
    \code{std\_ulogic}    & \code{logic}     & 1, unsigned  & 4 \\
    \code{std\_logic}     & \code{wire}      & 1, unsigned  & 4, resolved \\
    \code{boolean}        & ---              & use \code{bit} or \code{logic}; any non-zero is true & --- \\
    ---                   & \code{byte}      & 8, signed    & 2 \\
    ---                   & \code{shortint}  & 16, signed   & 2 \\
    \code{integer}        & \code{int}       & 32, signed   & 2 \\
    ---                   & \code{integer}   & 32, signed   & 4 \\
    ---                   & \code{longint}   & 64, signed   & 2 \\
    \code{natural}        & ---              & \code{int unsigned} is close but unconstrained & 2 \\
    \code{positive}       & ---              & no equivalent & --- \\
    \code{real}           & \code{real}      & 64-bit IEEE 754 & --- \\
    ---                   & \code{shortreal} & 32-bit IEEE 754 & --- \\
    \code{time}           & \code{time}      & 64, unsigned; see \cref{sec:rosetta:time} & 4 \\
    \code{character}      & \code{byte}      & 8, signed    & 2 \\
    \code{string}         & \code{string}    & dynamic in \SV{}, fixed length in \VHDL{} & --- \\
    \bottomrule
  \end{tabular}
\end{table}

Four entries deserve comment.

\textbf{\code{boolean}.} \VHDL{} has a distinct boolean type, and a condition
must be of that type: \code{if a = '1' then} is required because \code{a}
itself is not a boolean. \SV{} has no boolean type. A condition is any
expression; it is true if it is non-zero and, in a four-state type, known.
\code{if (a)} is therefore idiomatic and correct, and \code{if (a == 1'b1)} is
merely redundant. Where \VHDL{} forces you to be explicit, \SV{} does not, and
this is the single most common cosmetic change in a translated file.

\textbf{\code{int} versus \code{integer}.} \kw{int} is 32-bit signed two-state.
\kw{integer} is 32-bit signed four-state and exists only for Verilog
compatibility. \VHDL{}'s \code{integer} maps to \SV{}'s \kw{int}, not to
\SV{}'s \kw{integer}. Use \kw{int} for loop counters and for parameters. If you
need a loop variable that can hold \code{x}, you have a design problem, not a
type problem.

\textbf{\code{natural} and \code{positive}.} These are \VHDL{} subtypes with
range constraints, checked at run time. \SV{} has \code{int unsigned}, which
constrains the interpretation but not the range: nothing stops you assigning
\code{-1} and reading back \code{4294967295}. See the trap below.

\textbf{\code{string}.} A \VHDL{} string is a constrained array of
\code{character} and its length is part of its subtype, which is why
concatenating strings of different lengths is such a nuisance. A \SV{}
\kw{string} is a dynamically sized built-in type with methods
(\code{.len()}, \code{.substr()}, \code{.toupper()}) and ordinary \code{==}
comparison. This is one of the few places where \SV{} is unambiguously more
pleasant.

\begin{gotcha}{There are no range constraints and no run-time checks}
\begin{rosetta}
\begin{rleft}
subtype small_t is
  natural range 0 to 15;
signal s : small_t;

-- assigning 16 is a run-time
-- error: "value 16 out of
-- range 0 to 15"
s <= 16;
\end{rleft}
\begin{rright}
// no subtype constraints exist.
typedef int unsigned small_t;
small_t s;

// this is simply legal:
s = 16;
s = -1;      // reads as
             // 4294967295
\end{rright}
\end{rosetta}
Every \VHDL{} run-time range check you rely on disappears. A state variable
declared as \code{natural range 0 to 3} becomes \code{logic [1:0]}, which
wraps instead of failing, or \code{int}, which silently holds anything. The
replacement is explicit: an \code{assert} in the \RTL{}, an
\code{enum} for state variables, or an assertion in the testbench. Do not
assume the tool will tell you.
\end{gotcha}

\subsection{numeric\_std and arithmetic types}

\index{numeric\_std}\index{signed}\index{unsigned}

\begin{rosetta}
\begin{rleft}
use ieee.numeric_std.all;

signal u  : unsigned(7 downto 0);
signal s  : signed(7 downto 0);
signal v  : std_logic_vector(7 downto 0);
signal i  : integer;

u <= unsigned(v);
v <= std_logic_vector(u);
i <= to_integer(u);
u <= to_unsigned(i, 8);
s <= to_signed(i, 8);
u <= resize(u, 8);
u <= shift_left(u, 2);
s <= shift_right(s, 2);  -- arithmetic
u <= rotate_left(u, 3);
\end{rleft}
\begin{rright}
logic        [7:0] u;
logic signed [7:0] s;
logic        [7:0] v;
int                i;

assign u = v;         // no cast needed
assign v = u;         // no cast needed
assign i = u;         // zero-extends
assign u = i;         // truncates
assign s = i;         // truncates
// resize: implicit everywhere
assign u = u << 2;
assign s = s >>> 2;   // arithmetic
assign u = {u[4:0], u[7:5]}; // rotate
\end{rright}
\end{rosetta}

Look at the right-hand pane and count the conversion functions: there are none.
This is the central practical difference. In \VHDL{}, \code{unsigned},
\code{signed} and \code{std\_logic\_vector} are three distinct array types with
the same element type, and the type marks (\code{unsigned(v)}) are what let you
move between them. Every conversion is written down, and every one is checked.
In \SV{}, signedness is a property of the vector, not a separate type, so
\code{logic [7:0]} and \code{logic signed [7:0]} interoperate freely and the
compiler applies whatever conversion makes the expression legal.

The individual mappings:

\begin{itemize}
  \item \code{unsigned(v)} and \code{std\_logic\_vector(u)}: no operation at
        all in \SV{}. The bits are the bits.
  \item \code{to\_integer(u)}: implicit. Assigning a vector to an \kw{int}
        zero-extends if the vector is unsigned and sign-extends if it is
        signed. \code{\$signed(v)} and \code{\$unsigned(v)} force the
        interpretation where you need it.
  \item \code{to\_unsigned(i, N)} and \code{to\_signed(i, N)}: implicit
        truncation to the width of the target. There is no width argument
        because the target supplies it. If the value does not fit, the extra
        bits are dropped without a warning.
  \item \code{resize(u, N)}: implicit. This is convenient and it is the source
        of half the bugs in \cref{sec:rosetta:widths}.
  \item \code{shift\_left}/\code{shift\_right}: \code{<<}, \code{>>} and
        \code{>>>}. Note carefully that \code{>>>} is arithmetic only if its
        left operand is signed. On an unsigned operand it is identical to
        \code{>>}.
  \item \code{rotate\_left}/\code{rotate\_right}: no equivalent. Write the
        concatenation, or a small function.
\end{itemize}

\begin{tipbox}{A rotate function}
Since \SV{} has no rotate operator, put one in your package once:
\begin{svcode}
function automatic logic [31:0] rotl32 (logic [31:0] x, int unsigned n);
  return (x << n) | (x >> (32 - n));   // n in 1..31
endfunction
\end{svcode}
Note \kw{automatic}: see \cref{sec:rosetta:subprograms} for why this matters.
For a parameterised width, make the function part of a parameterised package or
give it a \code{parameter type}.
\end{tipbox}

\begin{notebox}{std\_logic\_arith and std\_logic\_unsigned}
You may have inherited code that writes \code{use ieee.std\_logic\_unsigned.all}
and then adds \code{std\_logic\_vector} values directly. Those packages are
Synopsys creations that were placed in the \code{ieee} library without ever
being IEEE standards. They overload the arithmetic operators on
\code{std\_logic\_vector} itself, which means the same vector cannot be treated
as signed in one expression and unsigned in another, and mixing them with
\code{numeric\_std} produces ambiguous overloads. Do not use them in new
\VHDL{}. There is an irony here worth noting: what those packages did to
\VHDL{}, by making a raw bit vector arithmetic and implicitly unsigned, is
exactly what \SV{} does by design. When you move to \SV{} you are permanently
in \code{std\_logic\_unsigned} mode, which is precisely why
\cref{sec:rosetta:widths} needs so much space.
\end{notebox}

% ===========================================================================
\section{Literals and aggregates}
\label{sec:rosetta:literals}
% ===========================================================================

\begin{rosetta}
\begin{rleft}
'0'            -- bit literal
'1'
'Z'
'X'
"1010"         -- 4-bit string literal
x"A5"          -- 8-bit hex
o"17"          -- 6-bit octal
b"1010_1010"   -- with separators
16#FF#         -- integer 255
2#1010#        -- integer 10
1_000_000      -- integer
1.5e-3         -- real
10 ns          -- physical
'a'            -- character
"hello"        -- string
\end{rleft}
\begin{rright}
1'b0           // sized bit literal
1'b1
1'bz
1'bx
4'b1010        // 4-bit binary
8'hA5          // 8-bit hex
6'o17          // 6-bit octal
8'b1010_1010   // with separators
255            // 32-bit int
10             // 32-bit int
1_000_000      // separators work
1.5e-3         // real
10ns           // time literal
"a"            // 8-bit value
"hello"        // string
\end{rright}
\end{rosetta}

\index{literal!sized}\index{literal!unsized}

\subsection{Sized and unsized}

A \VHDL{} literal is typed and, if it is an array literal, its length comes
from the context. \code{"1010"} is four bits because it has four characters;
\code{x"A5"} is eight bits because a hex digit is four bits. There is no way to
write a bit-vector literal whose length adapts.

A \SV{} literal has three parts: an optional width, a base, and the digits.
\code{8'hA5} is eight bits wide. \code{'hA5} without the width is unsized: it
takes its width from the context, and where there is no context it is at least
32 bits. A bare number such as \code{255} is an unsized decimal, which means a
32-bit signed \kw{int}.

Then there are the special unsized fill literals, added in \SV{}-2005, which
are the most useful literals in the language:

\begin{svcode}[caption={The fill literals.},label={lst:rosetta:fill}]
logic [31:0] a;
a = '0;    // 32 zeros   -- replaces (others => '0')
a = '1;    // 32 ones    -- replaces (others => '1')
a = 'x;    // 32 x
a = 'z;    // 32 z
a = 1'b1;  // 32'h0000_0001  -- NOT all ones
\end{svcode}

\begin{gotcha}{\code{'1} and \code{1'b1} are not the same thing}
\code{'1} fills every bit of the target. \code{1'b1} is a one-bit value that is
zero-extended to the target. In an eight-bit context the first is
\code{8'hFF} and the second is \code{8'h01}. The same distinction bites in
comparisons:
\begin{svcode}
logic [7:0] mask;
if (mask == '1)   ... // true when mask is 8'hFF   -- what you meant
if (mask == 1'b1) ... // true when mask is 8'h01   -- almost never
if (mask == 1)    ... // same as above: 1 is a 32-bit int
\end{svcode}
The \VHDL{} translation of \code{(others => '1')} is \code{'1}, always. Writing
\code{1'b1} where you meant \code{'1} produces code that is legal, silent, and
wrong.
\end{gotcha}

\begin{gotcha}{An unsized literal in a wide expression}
\code{'1} has no width of its own: it takes the width of its context. That is
usually what you want, but the context is not always what you think.
\begin{svcode}
logic [63:0] wide;
wide = {32'h0, '1};   // inside a concatenation every operand is
                      // self-determined. '1 has no size -> error or
                      // an implementation-defined width. Write 32'hFFFF_FFFF.
\end{svcode}
Unsized literals are illegal as concatenation operands for exactly this reason.
Some tools accept them and pick a width. Always give an explicit width inside
\code{\{\}}.
\end{gotcha}

\subsection{Aggregates and assignment patterns}

\begin{rosetta}
\begin{rleft}
constant ZERO : std_logic_vector(7 downto 0)
  := (others => '0');
constant ONE : std_logic_vector(7 downto 0)
  := (others => '1');
constant M : std_logic_vector(7 downto 0)
  := (3 downto 0 => '1', others => '0');
constant P : std_logic_vector(3 downto 0)
  := ('1', '0', '1', '0');
constant N : std_logic_vector(3 downto 0)
  := (3 => '1', 1 => '1', others => '0');

type pair_t is record
  x : integer; y : integer;
end record;
constant C : pair_t := (x => 1, y => 2);
\end{rleft}
\begin{rright}
localparam logic [7:0] ZERO = '0;

localparam logic [7:0] ONE = '1;

localparam logic [7:0] M = 8'h0F;
// or: {4'b0, {4{1'b1}}}

localparam logic [3:0] P = 4'b1010;

localparam logic [3:0] N = 4'b1010;


typedef struct packed {
  int x; int y;
} pair_t;
localparam pair_t C = '{x:1, y:2};
\end{rright}
\end{rosetta}

\index{aggregate}\index{assignment pattern}

For packed vectors the \VHDL{} aggregate almost always collapses to an ordinary
literal in \SV{}, because a packed vector \emph{is} a number. \code{(others =>
'0')} becomes \code{'0}. A named aggregate that sets specific bits becomes a
hexadecimal constant, a concatenation, or a part-select assignment:

\begin{svcode}[caption={Setting a field of a vector.},label={lst:rosetta:fieldset}]
logic [31:0] ctrl;
always_comb begin
  ctrl          = '0;      // (others => '0')
  ctrl[3:0]     = 4'hF;    // (3 downto 0 => '1', ...)
  ctrl[15:8]    = len_q;   // a named field
end
\end{svcode}

The \code{'\{...\}} form, the \emph{assignment pattern}, is the true equivalent
of a \VHDL{} aggregate, and it applies to unpacked arrays, structures and
unions. Note the leading apostrophe: \code{\{1,2,3\}} is a concatenation,
\code{'\{1,2,3\}} is an assignment pattern, and they mean entirely different
things.

\begin{svcode}[caption={Assignment patterns for composite types.},label={lst:rosetta:patterns}]
typedef struct packed { logic [7:0] a; logic [7:0] b; } two_t;
two_t   s = '{a: 8'h12, b: 8'h34};   // named, like VHDL's (a => ..., b => ...)
two_t   t = '{8'h12, 8'h34};         // positional

logic [7:0] mem [0:3] = '{8'h00, 8'h11, 8'h22, 8'h33};
logic [7:0] all9 [0:3] = '{default: 8'h99};   // like (others => x"99")
logic [7:0] two [0:3]  = '{2{8'hAA, 8'hBB}};  // replicated pattern
\end{svcode}

\code{'\{default: v\}} is the closest thing to \code{(others => v)} for an
unpacked array or a structure, and it works at any nesting depth.

\subsection{Concatenation and replication}

\begin{rosetta}
\begin{rleft}
y <= a & b & c;

-- replication needs a loop or an
-- aggregate:
y <= (others => a(7));
-- or
y <= a(7) & a(7) & a(7) & a(7);

-- sign extension by hand:
s16 <= (15 downto 8 => s8(7)) & s8;
-- or, with numeric_std:
s16 <= std_logic_vector(
         resize(signed(s8), 16));
\end{rleft}
\begin{rright}
assign y = {a, b, c};

// replication operator:
assign y = {4{a[7]}};

// nested:
assign y = {{4{a[7]}}, a[3:0]};

// sign extension:
assign s16 = {{8{s8[7]}}, s8};
// or, since s8 is signed:
assign s16 = $signed(s8);
\end{rright}
\end{rosetta}

\index{concatenation}\index{replication}

\code{\&} becomes \code{\{\}}. Note that \code{\&} in \SV{} is bitwise AND, so
mistranslating a concatenation gives you a legal expression with a completely
different meaning; the widths usually differ, but as \cref{sec:rosetta:widths}
explains, differing widths are not an error.

The replication operator \code{\{N\{expr\}\}} has no \VHDL{} counterpart short
of an aggregate or a loop, and it is used constantly, above all for sign
extension. The double brace is not a typo: the outer braces are the
concatenation, the inner ones delimit the replicated expression. \code{N} must
be a constant expression.

% ===========================================================================
\section{Operators, expressions and the width rules}
\label{sec:rosetta:widths}
% ===========================================================================

If you read only one section of this chapter, read this one. The operator
tables are boring and the width rules are not, and the width rules are where
\VHDL{} designers lose days.

\subsection{The operators}

\begin{table}[htbp]
  \centering
  \caption{Operators.}
  \label{tab:rosetta:operators}
  \small
  \begin{tabular}{@{}>{\raggedright\arraybackslash}p{0.235\linewidth}>{\raggedright\arraybackslash}p{0.215\linewidth}>{\raggedright\arraybackslash}p{0.47\linewidth}@{}}
    \toprule
    \VHDL{} & \SV{} & Note \\
    \midrule
    \code{and or xor xnor nand nor} & \code{\& | \textasciicircum{} \textasciitilde\textasciicircum{} \textasciitilde\& \textasciitilde|} & Bitwise. \\
    \code{not} & \code{\textasciitilde} & Bitwise. \\
    --- & \code{\&\& || !} & Logical (result is one bit). \VHDL{} uses the same words for both. \\
    \code{= /=} & \code{== !=} & \code{x} on either side gives \code{x}. \\
    --- & \code{=== !==} & Four-state identity: compares \code{x} and \code{z} literally. \\
    --- & \code{==? !=?} & Wildcard: \code{x}/\code{z} in the right operand match anything. \\
    \code{?= ?/=} (2008) & \code{==? !=?} & Nearest equivalent. \\
    \code{< <= > >=} & \code{< <= > >=} & Signedness comes from the operands in \SV{}. \\
    \code{+ - * /} \code{mod rem} & \code{+ - * / \%} & \code{\%} follows the sign of the dividend, like \code{rem}. \\
    \code{**} & \code{**} & \\
    \code{abs} & --- & No operator. Write \code{(x[N-1] ? -x : x)}. \\
    \code{sll srl sla sra rol ror} & \code{<< >> <<< >>>} & No rotate. \code{>>>} is arithmetic only on a signed operand. \\
    \code{\&} & \code{\{\}} & Concatenation. \\
    --- & \code{\{N\{x\}\}} & Replication. \\
    --- & \code{\&x |x \textasciicircum{}x} & Unary reduction over all bits. No \VHDL{}-2008 equivalent except \code{and\_reduce} etc. \\
    \code{when else} & \code{?:} & Conditional expression. \\
    \bottomrule
  \end{tabular}
\end{table}

Three rows deserve attention.

\textbf{Logical versus bitwise.} \VHDL{} uses \code{and} for both, and the type
system tells them apart: \code{and} on two booleans is logical, on two vectors
is bitwise. \SV{} has separate operators. \code{a \& b} on two eight-bit
vectors gives an eight-bit vector; \code{a \&\& b} gives one bit, true if both
operands are non-zero. Writing \code{if (a \& b)} where you meant \code{\&\&}
is legal and usually still works, which makes it a habit that survives until
the day it does not.

\textbf{The four-state comparisons.} \code{==} propagates unknowns: if either
operand contains an \code{x}, the result is \code{x}, and \code{if} treats
\code{x} as false. \code{===} compares the four-state values bit by bit and
always returns \code{0} or \code{1}, so \code{a === 8'hxx} is a way of asking
``is a exactly all-x''. \code{===} is not synthesisable and belongs in
testbenches only. \VHDL{} has no equivalent to \code{==} because
\code{std\_logic} equality already compares the nine values literally: \VHDL{}'s
\code{=} behaves like \SV{}'s \code{===}, not like \code{==}. This is a real
difference in behaviour, not just in spelling.

\textbf{Reduction operators.} \code{|v} is one bit, true if any bit of \code{v}
is set. \code{\&v} is true if all are. \code{\textasciicircum{}v} is the parity.
\VHDL{}-2008 added \code{or\_reduce}, \code{and\_reduce} and \code{xor\_reduce}
for the same purpose but the \SV{} form is one character.

\subsection{Width rules: the part that matters}

\index{width rules}\index{self-determined}\index{context-determined}

\VHDL{}'s rule is short. The operands of \code{+} on \code{unsigned} must be of
the same length or one of them must be an integer; the result has the length of
the longer operand; assigning it to a target of a different length is an error.
Every width is written down and every mismatch is a compile error.

\SV{} has no such rule. Instead, every operand of an expression is classified
as \emph{self-determined} or \emph{context-determined}. A self-determined
operand is evaluated at its own natural width. A context-determined operand is
first widened to the width of the whole expression, and the width of the whole
expression is the maximum of all its context-determined operands \emph{and,
for an assignment, of the left-hand side}. Operands are then zero-extended or
sign-extended to that width, the operation happens, and the result is
truncated to fit the target.

That is the whole rule, and it has three consequences you must internalise.

\begin{rosetta}
\begin{rleft}
signal a, b : unsigned(7 downto 0);
signal c    : unsigned(15 downto 0);

-- illegal: lengths differ
-- c <= a + b;

-- you must say what you mean:
c <= resize(a, 16) + resize(b, 16);
c <= "00000000" & (a + b);  -- 8-bit
                            -- add
\end{rleft}
\begin{rright}
logic  [7:0] a, b;
logic [15:0] c;

// legal. The 16-bit target makes
// the addition 16 bits wide, so
// a and b are zero-extended
// first and no carry is lost:
assign c = a + b;
\end{rright}
\end{rosetta}

\textbf{Consequence one: the target width propagates into the expression.} This
is the opposite of what most \VHDL{} designers expect, and it is usually
helpful. \code{c = a + b} with a sixteen-bit \code{c} really does compute a
sixteen-bit sum. The carry out of bit 7 is kept.

\textbf{Consequence two: the propagation stops at self-determined operands.}
The operands that are self-determined are the ones inside a concatenation or a
replication, the right-hand operand of a shift, the condition of \code{?:}, the
index of a part-select, an argument of a function or a task, and a port
connection. In every one of those places the expression is evaluated at its
own width and the surrounding context is ignored.

\begin{gotcha}{Where context stops: the carry you did not get}
\begin{svcode}
logic  [7:0] a, b;
logic [15:0] c;
logic        cout;

assign c    = a + b;             // 16-bit add. Carry kept. Fine.
assign c    = {8'h00, a + b};    // 8-bit add inside {}. Carry LOST.
assign cout = (a + b) >> 8;      // the whole expression is 8 bits wide, so
                                 // a+b wraps and >>8 is always 0. LOST.
assign {cout, c[7:0]} = a + b;   // 9-bit LHS -> 9-bit add. Carry kept.
\end{svcode}
The rule of thumb that protects you in every case, whatever the context, is to
widen an operand explicitly. \code{\{1'b0, a\} + b} is nine bits wide by
itself, regardless of where you put it, and it is exactly the \SV{} spelling of
\code{resize(a,9) + resize(b,9)}. Get into the habit and the width rules stop
being able to hurt you.
\end{gotcha}

\begin{tipbox}{Widen the operand, never trust the target}
\begin{svcode}
// Full adder result with explicit carry, safe anywhere:
assign sum9 = {1'b0, a} + {1'b0, b} + cin;

// Signed multiply-accumulate, safe anywhere:
assign acc = $signed({{8{a[7]}}, a}) * $signed({{8{b[7]}}, b});
\end{svcode}
Writing the extension by hand costs you a handful of characters and makes the
expression's width independent of where it appears. This is the single most
useful defensive habit in \SV{} \RTL{}.
\end{tipbox}

\textbf{Consequence three: truncation is silent.} Assigning a wide expression
to a narrow target drops the high bits with no message from most simulators.
\tool{Verilator} warns (\code{WIDTHTRUNC}, \code{WIDTHEXPAND}) and this is one
of the strongest arguments for lint-checking every file with it. Treat a
\tool{Verilator} width warning as an error.

\subsection{Signedness}

\index{signed!expression rules}

The signedness rule is one line, and it is the source of the second family of
disasters:

\begin{ruleofthumb}
If any operand of an expression is unsigned, the whole expression is unsigned.
\end{ruleofthumb}

Signedness is decided over the whole expression, before any values are
computed, and it is decided pessimistically. Two more facts complete the
picture. First, a concatenation is \emph{always} unsigned, whatever it
contains. Second, a part-select or bit-select is always unsigned, even of a
signed vector: \code{s[7:0]} is unsigned even if \code{s} is
\code{logic signed [7:0]}, though the whole-vector reference \code{s} is
signed.

\begin{gotcha}{One unsigned operand poisons the whole expression}
\begin{svcode}
logic signed [7:0] a;      // want -3
logic        [7:0] b;      //  want 5
logic signed [8:0] r;

assign a = -8'sd3;
assign b =  8'd5;
assign r = a + b;          // b is unsigned -> the WHOLE add is unsigned.
                           // a is zero-extended, not sign-extended:
                           // 9'd253 + 9'd5 = 9'd258, not 2.

assign r = a + $signed(b); // correct: 9'sd(-3) + 9'sd5 = 2

logic signed [7:0] x, y;
assign r = {x, 1'b0} + y;  // concatenation is unsigned -> whole add unsigned
\end{svcode}
The compiler says nothing. The waveform shows a plausible-looking number. This
is the bug that \code{numeric\_std} exists to prevent, and \SV{} gives you no
protection whatsoever: declare every signed quantity \code{logic signed}, and
put \code{\$signed()} around anything you mix with it.
\end{gotcha}

\begin{gotcha}{Unsigned comparison with a negative}
\begin{svcode}
logic [7:0] cnt;
if (cnt - 1 < 0)   ...     // NEVER true: unsigned result cannot be < 0
if (cnt < -1)      ...     // -1 is a 32-bit signed int, but cnt is unsigned,
                           // so the comparison is done unsigned:
                           // cnt < 32'hFFFF_FFFF -> always true
\end{svcode}
\VHDL{} is no better on the first of these. \code{numeric\_std} defines both
\code{"-"(UNSIGNED, NATURAL)} and \code{"<"(UNSIGNED, NATURAL)}, so
\code{cnt - 1 < 0} analyses cleanly and is likewise always false; only a lint
tool will tell you. \VHDL{} does save you from the second: the only matching
operator takes a \code{NATURAL} on the right, and \code{-1} is outside that
subtype, so \code{cnt < -1} is a bounds violation. An analyser will normally
warn about it, and it is a hard error when the expression is
evaluated. \SV{} accepts both, and the second
is worse than useless: it is a tautology that the tool will happily optimise
away, taking your intended logic with it.
\end{gotcha}

\subsection{The width of a comparison and of a ternary}

The result of a comparison is always one bit. Its \emph{operands}, however, are
context-determined with each other: they are widened to the maximum of the two,
and only then compared. So \code{a > b} with an eight-bit \code{a} and a
sixteen-bit \code{b} does a sixteen-bit comparison, which is what you want.
Beware the intermediate: \code{(a + b) > c} sizes \code{a + b} against
\code{c}, so if \code{c} is eight bits the addition wraps at eight bits before
the comparison happens.

The conditional operator is different again. In \code{cond ? x : y}, the
condition is self-determined, and \code{x} and \code{y} are context-determined
with the whole expression \emph{and with each other}. If \code{x} is signed and
\code{y} is unsigned, the result is unsigned, and both are treated as unsigned.

\begin{svcode}[caption={Widths in comparisons and conditionals.},label={lst:rosetta:cmpwidth}]
logic  [7:0] a, b;
logic [15:0] c;
logic [15:0] r;

if ((a + b) > c)  ...   // 16-bit add (c is 16 bits) then 16-bit compare
if ((a + b) > 8'd200) ...  // 8-bit add: wraps. 8'd200 is 8 bits.
if ((a + b) > 200)    ...  // 200 is a 32-bit int: 32-bit add. No wrap.

assign r = sel ? a : c;     // a is widened to 16 bits (context = r and c)
assign r = {8'h0, sel ? a : b};  // inside {}: self-determined, 8-bit result
\end{svcode}

The three \code{if} lines differ only in how the constant is written, and they
compute three different things. That is the width rule in one picture.

% ===========================================================================
\section{Arrays, records and enumerations}
\label{sec:rosetta:arrays}
% ===========================================================================

\subsection{Packed and unpacked}

\index{array!packed}\index{array!unpacked}

\SV{} splits the idea of an array in two, and this split has no \VHDL{}
counterpart at all. Dimensions written \emph{before} the object name are
\emph{packed}; dimensions written \emph{after} it are \emph{unpacked}.

\begin{rosetta}
\begin{rleft}
-- an array of 4 bytes
type mem_t is array (0 to 3) of
  std_logic_vector(7 downto 0);
signal m : mem_t;

-- element access
m(2)          -- one byte
m(2)(7)       -- one bit

-- VHDL makes no distinction
-- between "packed" and
-- "unpacked": the array is the
-- array.
\end{rleft}
\begin{rright}
// unpacked: 4 separate bytes
logic [7:0] u [0:3];

// packed: one 32-bit vector
// viewed as 4 bytes
logic [3:0][7:0] p;

// element access is identical
u[2]      p[2]      // one byte
u[2][7]   p[2][7]   // one bit

// but only p can do this:
p[19:4]   // arbitrary bit slice
p + 1     // arithmetic
\end{rright}
\end{rosetta}

\begin{figure}[htbp]
  \centering
  \begin{tikzpicture}[
      font=\footnotesize\sffamily,
      cell/.style={draw=codeframe,line width=0.7pt,minimum width=0.85cm,
                   minimum height=0.55cm,inner sep=1pt},
      grp/.style={draw=codekey,line width=1.1pt,rounded corners=1pt},
      lbl/.style={font=\scriptsize\sffamily\color{codenumber}}
    ]
    % --- packed -------------------------------------------------------------
    \node[lbl,anchor=west,font=\small\ttfamily\color{black}]
      at (0,1.35) {logic [3:0][7:0] p;};
    \node[lbl,anchor=west] at (0,0.95)
      {one contiguous 32-bit vector: \texttt{p[19:4]} and \texttt{p+1} are legal};
    \foreach \i in {0,...,31} {
      \pgfmathsetmacro{\x}{0.30*\i}
      \node[cell,minimum width=0.30cm,fill=svbg] at (\x+0.15,0.35) {};
    }
    \foreach \g/\name in {0/3,1/2,2/1,3/0} {
      \pgfmathsetmacro{\xa}{2.4*\g}
      \draw[grp] (\xa,0.06) rectangle (\xa+2.4,0.64);
      \node[lbl] at (\xa+1.2,-0.22) {\texttt{p[\name]}};
    }
    \node[lbl,anchor=west] at (0,-0.62) {bit 31};
    \node[lbl,anchor=east] at (9.6,-0.62) {bit 0};

    % --- unpacked -----------------------------------------------------------
    \node[lbl,anchor=west,font=\small\ttfamily\color{black}]
      at (0,-1.65) {logic [7:0] u [0:3];};
    \node[lbl,anchor=west] at (0,-2.05)
      {four separate 8-bit objects: no slice across them, no arithmetic};
    \foreach \g/\name in {0/0,1/1,2/2,3/3} {
      \pgfmathsetmacro{\xa}{2.4*\g + 0.35*\g}
      \foreach \i in {0,...,7} {
        \pgfmathsetmacro{\x}{\xa + 0.30*\i}
        \node[cell,minimum width=0.30cm,fill=vhdlbg] at (\x+0.15,-2.65) {};
      }
      \draw[grp,draw=accentalt] (\xa,-2.94) rectangle (\xa+2.4,-2.36);
      \node[lbl] at (\xa+1.2,-3.22) {\texttt{u[\name]}};
    }
  \end{tikzpicture}
  \caption{Packed and unpacked dimensions. A packed array is a bit vector with
           a grouping imposed on it; an unpacked array is a collection of
           separate objects. \VHDL{} has only the second kind, and gets the
           first through explicit conversion functions.}
  \label{fig:rosetta:packing}
\end{figure}

\cref{fig:rosetta:packing} is the picture to keep in mind. A packed array is a
single vector with a shape; it is stored contiguously, it can be sliced at any
bit boundary, it can be used in arithmetic, and it can be cast to and from any
other packed type of the same total width. An unpacked array is a collection of
separate objects; it is what you want for a memory, and it is what synthesis
tools recognise as a RAM.

\begin{table}[htbp]
  \centering
  \caption{What you may do with each kind.}
  \label{tab:rosetta:packing}
  \small
  \begin{tabular}{@{}lll@{}}
    \toprule
    Operation & Packed & Unpacked \\
    \midrule
    Whole-object assignment       & yes & yes (types must match exactly) \\
    Arithmetic and bitwise operators & yes & no \\
    Arbitrary bit slice           & yes & no \\
    Element and part select       & yes & yes (one element at a time) \\
    Module port                   & yes & yes \\
    Cast to another type          & yes & only through a bit-stream cast \\
    Element type may be unpacked  & no  & yes \\
    Inferred as a RAM             & no (it is one big register) & yes \\
    \code{foreach} iteration      & yes & yes \\
    \bottomrule
  \end{tabular}
\end{table}

The \VHDL{} translation rule is simple in practice. A \VHDL{} array of vectors
that you use as a memory becomes an unpacked array of packed vectors. A \VHDL{}
array of vectors that you slice, concatenate or convert to a flat vector
becomes a packed array.

\begin{svcode}[caption={Memories and buses.},label={lst:rosetta:memvsbus}]
// A 256 x 32 memory. Unpacked outer dimension: synthesis infers a RAM.
logic [31:0] ram [0:255];
always_ff @(posedge clk) if (we) ram[addr] <= wdata;

// An 8-lane bus you also want as one 256-bit word. Packed both ways.
logic [7:0][31:0] lanes;
assign flat_256 = lanes;              // legal: same total width
assign lanes[3] = something;          // per-lane access still works
\end{svcode}

\begin{gotcha}{Two dimensions, two meanings, one syntax}
\code{logic [3:0] x [0:7]} is eight four-bit words. \code{logic [7:0][3:0] y}
is one thirty-two-bit vector grouped into eight nibbles. \code{logic [3:0][7:0] z}
is one thirty-two-bit vector grouped into four bytes. They are different types
and the difference is which side of the name the dimension is on.

Worse, the index order reverses. For \code{logic [1:0][7:0] q}, \code{q[1]} is
the \emph{upper} byte, bits 15 down to 8; the leftmost packed dimension is the
outermost. For unpacked arrays the leftmost dimension is also the outermost.
When you have both, you read left to right through the packed dimensions and
then left to right through the unpacked ones, but the packed ones are written
before the name and the unpacked ones after it, so on the page the order is
reversed. Write the declaration, then write \code{\$bits()} of it in a quick
testbench and check.
\end{gotcha}

\subsection{Multidimensional arrays}

\begin{rosetta}
\begin{rleft}
-- a true 2-D array of scalars
type grid_t is array
  (0 to 3, 0 to 7) of std_logic;
signal g : grid_t;
-- indexed with a comma:
g(2, 5)

-- an array of arrays
type row_t is array (0 to 7)
  of std_logic;
type grid2_t is array (0 to 3)
  of row_t;
signal h : grid2_t;
h(2)(5)
\end{rleft}
\begin{rright}
// SystemVerilog has no comma-
// indexed arrays. Every multi-
// dimensional array is an array
// of arrays.
logic g [0:3][0:7];
g[2][5]

// same thing, packed:
logic [3:0][7:0] h;
h[2][5]

// slices of unpacked arrays:
logic row [0:7];
assign row = g[2];   // legal, but see below
\end{rright}
\end{rosetta}

\VHDL{}'s comma-indexed multidimensional array has no \SV{} equivalent, and it
loses nothing: the array-of-arrays form does everything the comma form does and
also lets you name and pass a row. \SV{} allows unpacked array slicing
(\code{g[2]} above yields a whole row) and whole-array assignment when the types
match exactly, which \VHDL{} also allows. Tool support for the slice in a
\emph{continuous} assignment is uneven: \tool{Verilator} 5.020 accepts
\code{assign row = g[2];}, while \tool{Icarus Verilog}~12 replies
\code{sorry: Array slices are not yet supported for continuous assignment}.
Inside an \code{always\_comb} block both are happy.

\subsection{Records and structures}

\begin{rosetta}
\begin{rleft}
type req_t is record
  valid : std_logic;
  addr  : std_logic_vector(31 downto 0);
  data  : std_logic_vector(31 downto 0);
end record req_t;

signal r : req_t;

r.valid <= '1';
r <= (valid => '1',
      addr  => a,
      data  => d);

-- a record cannot be treated as
-- a bit vector without a
-- conversion function you write
-- yourself.
\end{rleft}
\begin{rright}
typedef struct packed {
  logic        valid;
  logic [31:0] addr;
  logic [31:0] data;
} req_t;                  // 65 bits

req_t r;

r.valid = 1'b1;
r = '{valid: 1'b1,
      addr:  a,
      data:  d};

// a packed struct IS a bit vector
assign flat65 = r;
assign r      = flat65;
assign r[63:32] = a;   // legal too
\end{rright}
\end{rosetta}

\index{struct!packed}\index{record}

The correspondence is close, with one important addition. A \VHDL{}
\kw{record} is always what \SV{} calls unpacked: it is a collection of named
fields with no defined bit layout, and turning it into a vector requires a
conversion function you write. A \SV{} \code{struct packed} has a defined
layout: the first field occupies the most significant bits, the last field the
least significant, and the whole structure is a bit vector of the sum of the
field widths. That makes packed structures the natural way to describe a bus
transaction, a register field layout, or a pipeline stage payload, and it is
one of the genuinely good things \SV{} gives you.

Restrictions: a packed structure may contain only packed types, so no
unpacked arrays, no \kw{real}, no \kw{string}. An unpacked structure (write
\code{struct} without \kw{packed}) may contain anything, can cross a port, but
cannot be used in arithmetic or sliced.

Unions follow the same split. A \code{union packed} requires all members to be
the same width and gives you a reinterpreting view of the same bits, which is
the classic way to describe an instruction word. A plain \code{union} is
unpacked and needs \code{\$cast} or \code{tagged} discipline. \VHDL{} has no
union at all.

\begin{svcode}[caption={A packed union: two views of the same 32 bits.},label={lst:rosetta:union}]
typedef struct packed { logic [6:0] funct7; logic [4:0] rs2, rs1;
                        logic [2:0] funct3; logic [4:0] rd;
                        logic [6:0] opcode; } r_type_t;
typedef union packed {
  logic [31:0] raw;
  r_type_t     r;
} instr_t;

instr_t i;
assign i.raw    = fetch_data;
assign is_add   = (i.r.opcode == 7'h33) && (i.r.funct3 == 3'b000);
\end{svcode}

\subsection{Enumerations}

\begin{rosetta}
\begin{rleft}
type state_t is
  (IDLE, RUN, DONE);
signal s : state_t;

s <= IDLE;
if s = RUN then ... end if;

state_t'pos(RUN)    -- 1
state_t'val(1)      -- RUN
state_t'succ(IDLE)  -- RUN
state_t'pred(DONE)  -- RUN
state_t'image(s)    -- "run"
state_t'left        -- IDLE
state_t'right       -- DONE
\end{rleft}
\begin{rright}
typedef enum logic [1:0] {
  IDLE, RUN, DONE
} state_t;
state_t s;

s = IDLE;
if (s == RUN) ...

int'(RUN)      // 1
state_t'(1)    // RUN
s.next()       // successor of s
s.prev()       // predecessor
s.name()       // "RUN"
s.first()      // IDLE
s.last()       // DONE
s.num()        // 3
\end{rright}
\end{rosetta}

\index{enumeration}

The concepts match, the spelling does not. Two differences matter.

First, the base type. A \VHDL{} enumeration is its own type with an implicit
integer position, and synthesis picks the encoding. A \SV{} enumeration has a
base type, which defaults to \kw{int} if you do not give one. That default is
wrong for \RTL{}: a state variable of type \kw{int} is thirty-two bits, is
two-state, and starts at zero rather than \code{x}, so a missing reset is
invisible. Always write the base type:

\begin{ruleofthumb}
\code{typedef enum logic [N-1:0] \{...\} state\_t;} Never let an \RTL{}
enumeration default to \kw{int}.
\end{ruleofthumb}

Second, the strictness. \VHDL{} will not let you assign an integer to an
enumeration without \code{'val}. \SV{} will not either, in principle: assigning
an integer expression to an enumeration variable requires an explicit cast
\code{state\_t'(expr)}. This is one of the very few places where \SV{} is as
strict as \VHDL{}, and it is worth knowing because it surprises people who
expect \SV{} to convert everything. Assigning \emph{from} an enumeration to an
integer or a vector is implicit and needs no cast.

The strictness stops at the cast, though, and this is where \VHDL{} is still
ahead. \code{state\_t'(2'b11)} compiles and produces a variable holding a value
that has no name. \VHDL{}'s type simply has no such value, and \code{'val} with
an out-of-range position is a run-time error. If the source of the integer is
untrusted, use \code{\$cast}, which checks and reports failure rather than
storing a nameless value; \cref{sec:rosetta:conversion} gives the form.

Explicit encodings work in both:

\begin{svcode}[caption={A one-hot encoded state type.},label={lst:rosetta:onehot}]
typedef enum logic [3:0] {
  IDLE = 4'b0001,
  ADDR = 4'b0010,
  DATA = 4'b0100,
  RESP = 4'b1000
} state_e;
state_e state_q, state_d;
// $onehot(state_q) is then a useful assertion
\end{svcode}

The methods in the right-hand pane are the whole set, and they are worth
learning together because the \VHDL{} attributes you would reach for are spread
across three different mechanisms in \SV{}.

\code{.name()} returns the enumeration label as a \kw{string}, and it is the
direct replacement for the \VHDL{} attribute \code{'image}. The standard way
to print a state in a message is
\begin{svcode}
$display("state = %s at %0t", state_q.name(), $time);
\end{svcode}
It returns an empty string if the value does not correspond to a label, which
is itself a useful check, and it alone justifies giving a state machine an
enumerated type: every log message and every assertion failure prints the state
by name with no lookup table.

\code{.next()} and \code{.prev()} are \code{'succ} and \code{'pred} with one
important difference: they \emph{wrap round}. \code{DONE.next()} is
\code{IDLE}, where \code{state\_t'succ(DONE)} is an error in \VHDL{}. Both take
an optional count, so \code{s.next(2)} skips one member. The wrap-around makes
``advance until you reach \code{last()}'' an infinite loop if the exit test is
written carelessly; iterate over the members by count instead:

\begin{svcode}
state_t s = state_t'(0);
for (int unsigned i = 0; i < s.num(); i++) begin
  $display("%0d: %s", i, s.name());
  s = s.next();
end
\end{svcode}

\code{.num()} is the number of members and has no \VHDL{} attribute
counterpart; \code{.first()} and \code{.last()} are \code{'left} and
\code{'right}. Note that all six are methods called on a \emph{value}, not
attributes of a type, so you need an object of the type to call them, and
\code{state\_t::num()} is not legal.

% ===========================================================================
\section{Type conversion}
\label{sec:rosetta:conversion}
% ===========================================================================

\begin{rosetta}
\begin{rleft}
-- closely related array types:
u  <= unsigned(v);
v  <= std_logic_vector(u);

-- numeric conversions:
i  <= to_integer(u);
u  <= to_unsigned(i, 8);
s  <= to_signed(i, 8);

-- real and integer:
r  <= real(i);
i  <= integer(r);      -- rounds

-- enumeration:
p  <= state_t'pos(st);
st <= state_t'val(p);

-- anything else needs a function
-- you write yourself.
\end{rleft}
\begin{rright}
// related packed types: nothing
assign u = v;
assign v = u;

// numeric: implicit
assign i = u;
assign u = i;
assign s = i;

// real and integer: implicit,
// rounds to nearest
assign r = i;
assign i = r;

// enumeration:
assign p  = int'(st);
assign st = state_t'(p);

// static cast, any packed type:
assign q = my_t'(expr);
\end{rright}
\end{rosetta}

\index{cast!static}\index{cast!dynamic}

\SV{} has four conversion mechanisms, in increasing order of ceremony.

\textbf{Implicit conversion}, which happens whenever an expression of one type
is assigned to an object of another and the two are compatible. Between packed
types, ``compatible'' means ``both are collections of bits'', so almost everything
converts. Width mismatches are fixed by truncation or extension, signedness by
the rule in \cref{sec:rosetta:widths}. This is where the majority of \SV{} bugs
come from.

\textbf{The static cast} \code{type'(expr)}, which is the direct analogue of a
\VHDL{} type mark. It is checked at compile time. It works on any packed type,
on enumerations, and on \code{int}/\code{real} pairs. Note the apostrophe: the
syntax is the type name, an apostrophe, then the parenthesised expression, and
it looks disconcertingly like a \VHDL{} attribute.

\textbf{The size cast} \code{16'(expr)} and the sign casts
\code{\$signed(expr)}, \code{\$unsigned(expr)}, which change one property and
leave the rest alone. \code{16'(a+b)} is the cleanest way to force an
expression to a particular width without writing a concatenation, and it is
much clearer than relying on the target.

\textbf{The dynamic cast} \code{\$cast(dest, src)}, which checks at run time and
either assigns and returns 1, or leaves the destination alone and returns 0.
Use it when converting an integer to an enumeration where the value may not be
a valid label, and, in \refverificationbook, for downcasting class handles. It is a
function when you want the result and a task when you do not; used as a task,
a failed cast is an error.

\begin{svcode}[caption={The four mechanisms side by side.},label={lst:rosetta:casts}]
typedef enum logic [1:0] {A, B, C} e_t;
e_t          e;
logic [1:0]  v;
int          i;
logic [15:0] w;

assign v = e;              // implicit: enum to vector, always allowed
// e = v;                  // ILLEGAL: needs a cast
assign e = e_t'(v);        // static cast: no run-time check. v = 3 gives
                           //   an out-of-range enum, silently.
initial if (!$cast(e, v))  // dynamic cast: returns 0 if v is not A, B or C
  $error("bad state code %0d", v);
assign w = 16'(v);         // size cast: zero-extend to 16
assign i = $signed(v);     // sign cast: treat v's bits as signed
\end{svcode}

\begin{gotcha}{\SV{} converts where \VHDL{} refused to compile}
Every one of these produces a compile error in \VHDL{} and a silent conversion
in \SV{}:
\begin{svcode}
logic [7:0] a; logic [15:0] b; logic [3:0] c;
assign a = b;              // truncated to 8 bits
assign b = a;              // zero-extended to 16 bits
assign c = a + b;          // 16-bit add, then truncated to 4 bits
assign a = 1'b1;           // 8'h01
assign a = -1;             // 8'hFF
real r; assign r = a;      // vector to real
int  i; assign i = r;      // real to int, rounded
\end{svcode}
None of them is a mistake by itself. All of them are mistakes when they were
not intended, and nothing in the language tells you which is which. The
compensating discipline is: run \tool{Verilator} with \code{-Wall} over every
file, and treat \code{WIDTHTRUNC} and \code{WIDTHEXPAND} as errors unless you
have deliberately written a cast to say the conversion is wanted.
\end{gotcha}

\begin{notebox}{Bit-stream casts}
\SV{} also supports \emph{bit-stream} casts, which flatten an arbitrary
aggregate, including unpacked arrays and unpacked structures, into a bit vector
and back:
\begin{svcode}
typedef struct { logic [7:0] a; logic [7:0] b [0:3]; } msg_t;   // unpacked
msg_t        m;
logic [39:0] flat;
assign flat = {>>{m}};          // pack, left to right
assign m    = {>>{flat}};       // unpack
\end{svcode}
This has no \VHDL{} equivalent short of writing the serialisation by hand. It
is not synthesisable and it is meant for testbenches: it is the standard way to
turn a transaction object into a byte stream for a scoreboard or a file. The
\code{\{<<\{...\}\}} form streams right to left, which is how you reverse a bit
or byte order in one expression.
\end{notebox}

Type conversion is the natural place to stop, because it is where the static
half of the language runs out. Everything this chapter has described is
decided before the simulation starts: a width, a signedness, an encoding, a
layout in bits. None of it says when anything happens.

That is the subject of \cref{ch:rosetta2}, and the two halves meet almost
immediately. A continuous assignment applies the width rules of
\cref{sec:rosetta:widths} on every input change; an \code{always\_ff} decides
whether a truncation happens once per clock or never; the packed structure of
\cref{sec:rosetta:arrays} is what crosses a port in an instantiation. Keep the
rules below to hand while you read it, because the next chapter assumes them
rather than restating them.

% ===========================================================================
\section{Checklist}
\label{sec:rosetta:checklist}
% ===========================================================================

Operational takeaways for the static half of the language. The behavioural
ones are at the end of \cref{ch:rosetta2}, and the complete two-part mapping
table is with them.

\begin{itemize}
  \item \textbf{Every file starts} with \tick{timescale 1ns/1ps} and
        \tick{default\_nettype none}, and ends with
        \tick{default\_nettype wire}.
  \item \textbf{Ports:} \code{input wire}, \code{output logic},
        \code{inout wire}. ANSI style only. Connect every port explicitly;
        \SV{} has no port defaults and an unconnected input is \code{z}.
  \item \textbf{Generics become parameters}, architecture constants become
        \kw{localparam}. Never write \code{defparam}.
  \item \textbf{\code{logic} for everything} you drive from one place;
        \code{wire} only where you genuinely need resolution.
  \item \textbf{Declare signedness once}, in the declaration, and never mix a
        signed and an unsigned operand in one expression without a cast.
  \item \textbf{Widen before you add.} \code{\{1$^\prime$b0, a\} + b} for the
        carry; assigning to a wider target does not create one.
  \item \textbf{\code{>>>} on signed values}, always. Check every shift.
  \item \textbf{\code{$^\prime$0} and \code{$^\prime$1}} for context-width
        fills; a bare \code{0} is 32 bits and will surprise you inside a
        concatenation.
  \item \textbf{A packed struct is a bit vector} with names. Use it wherever
        you would have written a \VHDL{} record on a port.
  \item \textbf{Give every enumeration a base type} and, in \RTL{}, explicit
        encodings.
  \item \textbf{Lint every file} with \code{verilator -{}-lint-only -Wall}.
        It is the type checker \SV{} does not have.
\end{itemize}
